\chapter{Decisiones de diseño}
\label{cap:decisiones}

El código dice \emph{qué} hace el programa; este capítulo dice \emph{por qué} es así y qué se
probó y se descartó. Cada decisión que cruza áreas, o que fijó el propietario, está en un
documento corto llamado ADR (\emph{Architecture Decision Record}) en
\archivo{docs/decisions/}. Aquí tienes las dieciocho agrupadas por tema, más las preferencias
del propietario que no llegaron a ADR.

\begin{idea}[Por qué existen las ADR]
Dentro de seis meses, alguien (o el mismo agente tras compactar su contexto, que no recuerda
nada) verá algo raro, como el prefijo \cod{lavf://} o un temporizador a 10 Hz, y querrá
«arreglarlo». Si el motivo no está escrito, lo romperá. La regla del proyecto es:
\emph{no se reabre una ADR sin un motivo nuevo}.
\end{idea}

\section{Las prioridades que mandan}

Antes de las decisiones sueltas, la que las ordena. El propietario fijó cuatro prioridades,
\emph{en este orden}:

\begin{enumerate}
  \item Consumo bajo de memoria.
  \item Respuesta instantánea.
  \item Personalización total (vistas, disposición, skins).
  \item Estética cuidada.
\end{enumerate}

Ante un conflicto gana la de arriba, \emph{y se le explica}. El objetivo declarado es
sustituir a Pear Desktop y a una pestaña de Chrome con YouTube Music: «lo que no consuma
menos que ellos no vale». Casi toda decisión de este capítulo es esa frase aplicada.

\section{La base: pila y procesos}

\subsection{ADR-01: Qt/QML para la interfaz, Rust para el núcleo}

\begin{description}[style=nextline, leftmargin=0pt]
  \item[Se eligió] Qt 6/QML con un núcleo en Rust unido con cxx-qt. Los objetos que ve QML son
  los cinco singletons de Rust.
  \item[Se descartó] \emph{Electron o Tauri}: un WebView solo ya cuesta más que la app
  entera, y es exactamente lo que se quiere sustituir. \emph{GTK o interfaces nativas}: dos
  interfaces que mantener y sin el 3D ni el vidrio. \emph{Núcleo en C++}: posible, pero Rust
  da hilos y red más seguros.
  \item[Consecuencia] Toda la interfaz pasa por QML y sus trampas; cada archivo QML y cada
  puente hay que registrarlo en \archivo{build.rs}.
\end{description}

\subsection{ADR-02: mpv en un proceso aparte}

\begin{description}[style=nextline, leftmargin=0pt]
  \item[Se eligió] Un proceso \cod{mpv --idle} sin vídeo, controlado por su socket IPC JSON.
  \item[Se descartó] \emph{libmpv embebido}: otra dependencia nativa en dos sistemas, y un
  fallo de mpv tumbaría la app (queda pendiente medir si ahorra memoria). \emph{QtMultimedia}:
  no lee bien las URL de googlevideo ni Opus en todos los sistemas.
  \item[Por qué es buena idea] Si mpv se cuelga, se reinicia sin cerrar la app; ocupa 20 MB
  en reposo y unos 75 MB sonando; y su propio gancho de yt-dlp es el plan B de la resolución.
\end{description}

\subsection{ADR-07: la ventana de inicio de sesión, en otro proceso}

\begin{description}[style=nextline, leftmargin=0pt]
  \item[Se eligió] En macOS, el mismo ejecutable relanzado con \cod{--login} (WebKit del
  sistema, sin arrancar Qt); en Linux, un binario aparte (\cod{musicplayer-login}, Qt WebEngine,
  opcional).
  \item[Se descartó] \emph{Qt WebEngine dentro del reproductor}: cargaría Chromium (cientos
  de MB) para usarlo una vez. \emph{Disfrazar el agente de usuario de Chrome}: Google lo
  bloquea porque contradice las cabeceras \cod{Sec-CH-UA}. \emph{Solo importar del
  navegador}: en macOS exige «Acceso total al disco».
  \item[Lección general] La memoria es la primera prioridad: lo que se necesita una sola vez
  vive en un proceso que desaparece.
\end{description}

\section{Obtener el audio}

\subsection{ADR-03: vía rápida y yt-dlp de respaldo}

\begin{description}[style=nextline, leftmargin=0pt]
  \item[Se eligió] Pedir la URL a \cod{/player} con el cliente de visionOS (0,2--0,4 s). Si
  falla, yt-dlp. Las pistas que YouTube reserva a cuentas van directas a yt-dlp con las
  cookies.
  \item[Se descartó] \emph{Solo yt-dlp}: 3,5--4,2 s por pista. \emph{El gancho de yt-dlp de
  mpv}: 8,2--9,5 s. \emph{Otros clientes de InnerTube}: piden descifrado de JavaScript o
  tokens PO.
  \item[Consecuencia] La vía rápida sirve unas 9 de cada 10 pistas. Cuando YouTube cambie algo,
  lo primero que falle será ella: por eso yt-dlp debe seguir funcionando solo.
\end{description}

\subsection{ADR-05: precargar}

Las URL se resuelven por adelantado: el foco detenido 250 ms, el ratón sobre una carátula, las
tres pistas siguientes, las seis primeras de un álbum abierto y las cuatro primeras de una
sección. Con precarga el clic suena en 64--254 ms; sin ella, en 0,4--0,6 s.
\emph{Se descartó} precargar todo lo visible (decenas de peticiones por pantalla y riesgo de
que YouTube bloquee la IP). \cod{hint} atiende solo la \emph{última} petición.

\subsection{ADR-08: yt-dlp residente}

Arrancar Python en cada pista cuesta casi un segundo. Un yt-dlp que se queda cargado hace que
las pistas que lo necesitan suenen ~1 s antes (2,2--2,5 s frente a 3,5--4,2 s). Pero ocupa
unos 120 MB, así que se cierra solo a los 60 s sin uso y se puede desactivar. \emph{Se
descartó} dejarlo siempre vivo.

\section{Los datos}

\subsection{ADR-04: JSON en una propiedad \texttt{QString}}

\emph{Se eligió} publicar las listas como JSON en una propiedad de texto, que
\archivo{Catalog.qml} parsea una vez por cambio. \emph{Se descartó} \cod{QAbstractListModel}
(mucho código de puente por cada dato nuevo para un volumen pequeño) y \cod{QVariantList} (no
ahorra nada medible). El límite: 60 elementos por zona. Si un día hubiera listas de miles de
filas que cambian sueltas, habría que medir y quizá pasar esa lista a un modelo.

\subsection{ADR-06: las canciones de un álbum, de la página \texttt{VL<id>}}

La cola \cod{next} parecía lo natural, pero trae junto a cada canción su versión en vídeo, se
corta en unas 50 y a veces mezcla canciones ajenas. La página de la lista (\cod{VL} más el
identificador) trae solo la lista y en orden. Una lista de 500 pistas tarda ~3 s: se muestra
lo conocido al instante y las canciones después.

\subsection{ADR-09: caché de secciones en disco}

\emph{Esperar a la red} dejaba la ventana vacía medio segundo o más. Ahora cada sección se
guarda y se muestra al arrancar mientras se pide la nueva: a los 300 ms ya hay contenido y
carátulas. La contrapartida es que durante medio segundo lo que ves puede ser de la vez
anterior; y que cerrar sesión debe borrar las copias de la cuenta.

\subsection{ADR-12: registrar las escuchas}

\begin{description}[style=nextline, leftmargin=0pt]
  \item[Se eligió] Enviar los mismos avisos que el reproductor de YouTube Music
  (\cod{videostatsPlaybackUrl} al empezar y \cod{videostatsWatchtimeUrl} con los tramos), con
  \cod{c=WEB_REMIX}. Activado por defecto.
  \item[Se descartó] \emph{Confiar en la petición del audio} (no cuenta como reproducción) y
  \cod{yt-dlp --mark-watched} (lo apunta en el historial de YouTube, no en el de YouTube
  Music, y no envía el tiempo escuchado).
  \item[Consecuencia] YouTube aplica su límite de un dispositivo a la vez por cuenta. Nada
  espera a estos avisos; un fallo solo va a stderr.
\end{description}

\subsection{ADR-18: letra de varias fuentes}

\begin{description}[style=nextline, leftmargin=0pt]
  \item[Se eligió] Pedir LRCLIB, NetEase y Kugou \emph{a la vez} y publicar todas las que
  respondan, ordenadas de mejor a peor. La interfaz elige cuál mostrar.
  \item[Por qué] Solo con LRCLIB faltaba mucho anime y casi nunca había tiempos por palabra:
  el barrido se inventaba (75 ms por letra) y no encajaba con la voz.
  \item[Se descartó] \emph{Musixmatch} (token que se saca imitando su aplicación, captchas),
  \emph{QQ Music} (3DES modificado, catálogo ya cubierto), \emph{Apple Music y Spotify}
  (exigen la cuenta del usuario), \emph{alinear la letra con la voz en local con un modelo}
  (cientos de MB de memoria) y \emph{elegir la fuente en Rust} (obligaba a pasar las
  preferencias por canción al hilo de la letra; en QML cambiar de fuente es instantáneo porque
  todas ya están en \cod{Catalog.lyricsSources}).
\end{description}

\section{La interfaz}

\subsection{ADR-10: el mini reproductor oculta la ventana grande}

Al abrir el mini, la ventana grande se \emph{oculta} (no se minimiza) y \archivo{window.cpp}
libera su memoria gráfica: 213 $\rightarrow$ 148 MB. \emph{Se descartó} minimizar (en macOS Qt
conserva los búferes) y meter el mini en la misma ventana (no libera nada). El precio: las
\cod{Shape} con \cod{CurveRenderer} no se redibujan solas, y por eso existe
\cod{Theme.renderEpoch}.

\subsection{ADR-11: skins en JSON, opciones como mapa}

Una skin es un archivo JSON de colores; una opción es una pareja clave--valor en un mapa
genérico. \emph{Se descartó} una propiedad de puente por opción (cada una tocaría Rust, el
puente y QML) y las skins con QML propio (aplazadas; los colores cubren casi todo). El precio:
las claves no tienen tipo, así que cada una se documenta en la ficha de su área.

\subsection{ADR-13 y ADR-18: la letra y su coste}

\begin{itemize}
  \item \textbf{Un solo gestor de clics} para toda la vista de la letra, calculado con
  \cod{childAt}; un gestor por línea (o con «hover») costaba ~25 \% de CPU incluso sin
  animaciones.
  \item \textbf{Tres calidades} que cambian los fotogramas por segundo.
  \item \emph{Descartado}: \cod{Text.CurveRendering} (+95 MB para una nitidez que apenas se
  nota) y el brillo con \cod{layer.effect} (el texto se veía borroso).
\end{itemize}

\subsection{ADR-14: Lomos con matrices propias}

\emph{Se descartó} Qt Quick 3D (un motor 3D entero para cajas planas) y la \cod{Rotation} de
Qt por cara (el eje y el lomo visible cambiaban de lado de golpe al cruzar el centro). Cada
cara es un \cod{Item} plano con una \cod{Matrix4x4} calculada, y ángulo, escala y sitio son
funciones continuas de una posición decimal. El reflejo y el halo cuestan unos 7 MB.

\subsection{ADR-17: el fondo «Fluido»}

Un \cod{ShaderEffect} con un sombreador propio, precompilado con \cod{qsb} y guardado en el
repositorio junto a su \cod{.frag}. La carátula llega diminuta (32 px) y el desenfoque se
completa en el propio sombreador, dibujado a media resolución lógica. \emph{Se descartó}
copiar el sombreador de Better Lyrics Shaders (es GPL-3.0: solo se toma la técnica, tal como
la describe Kawarp, con licencia MIT), compilar con \cod{qsb} en cada \cod{cargo build} (obligaría a tener \cod{qsb}
instalado) y el desenfoque Kawase en varias texturas (con 32 px no hace falta). Quieto cuesta
0,08 s de CPU en 25 s y 196 MB (como sin él).

\section{Plugins y documentación}

\subsection{ADR-15: plugins en un hilo propio}

\cod{Player} avisa con \cod{plugins::emit} y un hilo reparte a Discord y Last.fm. \emph{Se
descartó} llamarlos desde \cod{Player}: la red de un plugin lento frenaría el arranque de una
pista o la interfaz. Los controles del sistema \emph{no} pasan por ahí porque en macOS deben
ir en el hilo principal.

\subsection{ADR-16: documentación por áreas}

Normas en \archivo{AGENTES.md}; una ficha por área; ADR; \emph{runbooks}; trampas; historial; y
\cod{scripts/docs.py} que lo comprueba (sin dependencias: Python 3.9 o más). \emph{Se
descartó} mover el código a una carpeta por funcionalidad (rompe las rutas planas que esperan
\archivo{build.rs} y QML) y un solo \archivo{AGENTS.md} largo (crecía con cada tarea y se leía
entero para cualquier cambio). El motivo es mantener el proyecto con agentes que pierden la
memoria al compactar el contexto.

\section{Preferencias del propietario que no son ADR}

En \archivo{docs/AGENTES.md} §11 hay una lista de preferencias vigentes que «no se cambian sin
preguntar». Son decisiones de producto más que de arquitectura. Las principales:

\begin{itemize}
  \item \textbf{Clic y que suene}: el arranque de una pista no espera a nada que no sea el
  audio.
  \item \textbf{Ventana de vidrio con navegación superior}; el panel de letra y cola, oculto al
  arrancar.
  \item \textbf{Vistas por sección}: Inicio, Explorar, Me gusta y Búsqueda en Muro;
  Playlists y Álbumes en Lomos; cualquier sección puede usar cualquier vista.
  \item \textbf{Diseño con alternativas}: mockups y decisiones visuales con dos a cuatro
  opciones a elegir, no una sola propuesta.
  \item \textbf{«Sonando» en dos disposiciones} (Grande y Clásica) y la carátula a máxima
  calidad por defecto.
  \item \textbf{La mascota participa}: pincharla da «Me gusta» si su \cod{mascota.json} lo pide.
  \item \textbf{Letra centrada} en horizontal y vertical; el retraso, un pie discreto.
  \item \textbf{Iniciar sesión con una ventana propia}, sin disfrazar el agente de usuario.
\end{itemize}

\begin{porque}[¿Cuándo escribo una ADR?]
Cuando una decisión cruza áreas, o la fija el propietario, o descartaste una alternativa
razonable y alguien podría querer reabrirla. El procedimiento es
\archivo{docs/runbooks/registrar-decision.md}: copias la plantilla, rellenas \emph{Decisión,
Alternativas descartadas, Motivo y Consecuencias} (con cifras y fecha cuando las haya) y
ejecutas \cod{python3 scripts/docs.py index}.
\end{porque}

\begin{ejercicio}
Elige una ADR y busca en el código dos sitios donde se note su \emph{consecuencia} (por
ejemplo, ADR-09: \cod{library::cache} y la llamada a \cod{forget_account}). Si quitaras la
decisión, ¿qué se rompería?
\end{ejercicio}
